\section{Main result: 200 paired episodes in clear weather}\label{sec:headline}

\begin{plainwords}
Seven drivers (and two rule-based teachers) drove the same 200 scenarios. The driver we propose succeeds in 92.5\% of them. The two baseline drivers that are committed on GitHub succeed in 85.0\% and 74.5\%. The proposed driver is statistically as good as the teacher it learned from.
\end{plainwords}
\begin{cautionbox}[Compared with GitHub]
The GitHub README states 90\% success (9 of 10 episodes) for the traffic-trained baseline. On 200 episodes the same checkpoint scores 74.5\% (interval 0.680 to 0.800); the README figure comes from ten episodes with the full camera pipeline, whose interval (0.60 to 0.98) is compatible with it. Table~\ref{tab:githubbaseline} gives the comparison.
\end{cautionbox}


\paragraph{Protocol.} Seeds 5000--5199, sparse traffic (density 0.05), clear weather, horizon 800 steps; the same 200 scenarios for every model; script \pth{bev_rl/adaptive/evaluate_final_200seed_sparse.py}; per-episode results in \pth{bev_rl/adaptive/results/final_200seed/}.
No model was trained, validated or selected on these seeds. Six models are compared: the two rule-based teachers (which see the exact simulator state), the proposed model, the caution-study control, and the two behaviour-cloning baselines.

\explainedfigure{fig_headline_success_clear_weather_200_episodes}{0.97}{Success rate of the driving models on the same 200 test episodes}{Each bar is the share of 200 identical test episodes (scenario seeds 5000 to 5199, clear weather, sparse traffic) in which a model reached the destination without a collision and without leaving the road. The black line is the 95\% Wilson confidence interval, the range of success rates that is statistically compatible with the observed count. Grey bars are rule-based teachers that see the exact simulator state; orange bars are the two baseline models supplied by the teammates; blue bars are models trained in this project.}{The Proposed Bird's-Eye-View Driving Model succeeds in 92.5\% of the episodes. The two baseline models succeed in 85.0\% (trained without traffic) and 74.5\% (trained with traffic). The proposed model is about as good as the Collision-Aware Teacher (95.0\%) whose behaviour it learned; the intervals overlap, so the two cannot be told apart. The Image-Only Driving Model, which reads camera pixels instead of the bird's-eye-view grid, reaches 0.130 (interval 0.090 to 0.184), far below every model that uses the grid.}{Training a model by letting it drive and asking a teacher what it should have done (dataset aggregation) produces a model that matches its teacher and clearly beats few-demonstration behaviour cloning. This figure does not say \emph{which part} of the proposed design (the training method or the map overlay) causes the gain; that needs the additional evaluations listed in the paper-directions section.}{fig:headline}


\begin{table}[H]\centering\small
\caption{Result table for the 200 clear-weather episodes (95\% Wilson intervals; ``Left road'' includes leaving the planned route; energy is the estimate of Section~\ref{sec:setup}).}\label{tab:headline}
\resizebox{\textwidth}{!}{\begin{tabular}{@{}lrrrrrrr@{}}\toprule
Model & Success & 95\% interval & Collision & Left road & Timeout & Speed (m/s) & Energy (Wh per km)\\\midrule
Collision-Aware Teacher (rule-based; sees exact simulator state) & 0.950 & [0.910, 0.973] & 0.010 & 0.040 & 0.000 & 7.84 & 101.3\\
Proposed Bird's-Eye-View Driving Model (map overlay + 3 vehicle-state inputs) & 0.925 & [0.880, 0.954] & 0.060 & 0.010 & 0.005 & 7.87 & 96.2\\
Caution-Study Control Model (caution input fixed to zero) & 0.895 & [0.845, 0.930] & 0.105 & 0.000 & 0.000 & 7.98 & 80.4\\
Stock Intelligent Driver Model (rule-based; sees exact simulator state) & 0.885 & [0.833, 0.922] & 0.080 & 0.035 & 0.000 & 7.90 & 97.9\\
Baseline Driving Model Trained Without Traffic & 0.850 & [0.794, 0.893] & 0.125 & 0.025 & 0.000 & 8.09 & 85.9\\
Baseline Driving Model Trained With Traffic & 0.745 & [0.680, 0.800] & 0.120 & 0.135 & 0.000 & 8.07 & 90.7\\
Image-Only Driving Model (reads front-camera pixels) & 0.130 & [0.090, 0.184] & 0.200 & 0.605 & 0.065 & 9.20 & 276.6\\
\bottomrule\end{tabular}
}
\end{table}

\subsection{The Image-Only Driving Model on the same 200 episodes}
The Image-Only Driving Model (front-camera picture of 128 by 96 pixels as its only input; Section~\ref{sec:arch}) was evaluated on the same 200 scenarios with the same metrics (script \pth{bev_rl/adaptive/vanilla_camera_cnn.py eval}; per-episode results in \pth{bev_rl/adaptive/results/vanilla_camera/}). It differs from the other models in the harness only in what it reads: it is driven by a rendered 320 by 240 camera image instead of the grid, so no detector or degradation model is involved (clear weather).
Result: success \ImOnSuccess{} (95\% interval \ImOnLow{} to \ImOnHigh{}); collisions \ImOnCollision{}, left the road \ImOnOffroad{}, timeouts \ImOnTimeout{}; mean speed \ImOnSpeed{} metres per second; estimated energy \ImOnEnergy{} watt-hours per kilometre. It is included in the figures and tables of this section.
\begin{implicationbox}
A driver that reads the picture directly is far below every driver that reads the grid. The grid is a compact, structured summary produced by the object detector and the stereo depth; with the amount of training here, a network cannot recover that summary from raw pixels. Whether more data would close the gap is not tested.
\end{implicationbox}

\subsection{What the numbers say}
\begin{itemize}
\item \textbf{The proposed model is clearly better than both baselines.} 0.925 against 0.850 (no-traffic baseline) and 0.745 (sparse-traffic baseline). The paired differences are $+0.075$ (95\% interval $[+0.035,+0.120]$, exact McNemar $p=0.0007$) and $+0.180$ ($[+0.125,+0.240]$, $p<0.0001$); see Figure~\ref{fig:paired} and Table~\ref{tab:paired}.
  The baselines' failures are of a different kind: the sparse-traffic baseline leaves the road in 13.5\% of the episodes (proposed model: 1.0\%), a sign of lateral drift, and the no-traffic baseline collides in 12.5\%.
\item \textbf{The proposed model is statistically indistinguishable from the collision-aware teacher} it is trained to imitate (difference $-0.025$, interval $[-0.070,+0.020]$, $p=0.40$), and $+0.040$ better than the stock driver model (interval $[0.000,+0.080]$, $p=0.077$, i.e.\ not significant at the 5\% level).
\item \textbf{The caution-study control (0.895) is not significantly different from the proposed model} ($+0.030$, $p=0.18$). Its apparent 0.975 on 40 episodes (Section~\ref{sec:weather}) was a selection effect of a small sample: the best of several candidates evaluated on 40 episodes regresses towards its true value when re-evaluated on 200.
\item \textbf{Energy.} The proposed model uses an estimated 96.2\,watt-hours per kilometre at 7.87\,m/s, similar to the stock driver model (97.9\,watt-hours per kilometre) and the collision-aware teacher (101.3). The two behaviour-cloning baselines use less (85.9 and 90.7\,watt-hours per kilometre) \emph{and drive faster} (8.09 and 8.07\,m/s) because they are less careful. The caution-study control is the cheapest model (80.4\,watt-hours per kilometre) at 7.98\,m/s with 89.5\% success;
  the reason for its lower consumption has not been investigated (Figure~\ref{fig:energy}).
\end{itemize}

\explainedfigure{fig_outcome_breakdown_200_episodes}{0.97}{How the 200 test episodes end, for each model}{For every model, the 200 episodes are split into the four possible outcomes: success (green), collision (red), left the road (yellow) and timeout (grey). Each bar sums to one.}{The proposed model fails almost only by collision (6.0\% collisions, 1.0\% left the road, 0.5\% timeouts). The baseline trained with traffic leaves the road in 13.5\% of the episodes and collides in 12.0\%; the baseline trained without traffic collides in 12.5\%. The Image-Only Driving Model ends 0.200 of the episodes in a collision, 0.605 by leaving the road and 0.065 by timeout.}{The baselines fail in a way that points to \emph{drift}: with only five demonstrations they never saw how to recover from a small steering error, so they slide off the road. Iterative training removes that failure type (1.0\% left the road). What remains for the proposed model is collisions with other vehicles, which points to what the model can see rather than to how it was trained.}{fig:outcomes}


\explainedfigure{fig_paired_difference_headline_versus_others}{0.97}{Paired difference in success between the proposed model and each other model}{For each comparison the difference (proposed model minus the other model) is computed on the same 200 episodes. The point is the difference; the black line is the 95\% bootstrap interval obtained by resampling episodes 5000 times. A line entirely to the right of the zero line means the proposed model is better with statistical support.}{The proposed model is better than the baseline trained without traffic by 0.075 (interval $+0.035$ to $+0.120$) and better than the baseline trained with traffic by 0.180 (interval $+0.125$ to $+0.240$). Its differences to the Collision-Aware Teacher ($-0.025$), the caution-study control ($+0.030$) and the stock driver model ($+0.040$) have intervals that include or touch zero. The Image-Only Driving Model is far below it.}{The claim that the proposed model beats the two baselines is statistically supported. No claim is made that it beats the teachers or the control. Using paired comparisons (same scenarios for every model) is what makes a difference of 0.075 detectable with 200 episodes.}{fig:paired}


\begin{table}[H]\centering\small
\caption{Paired comparison of the proposed model with each other model. ``Only proposed model'' counts the seeds on which the proposed model succeeded and the other model did not; the exact McNemar test uses only these discordant seeds.}\label{tab:paired}
\resizebox{\textwidth}{!}{\begin{tabular}{lrrrrr}\toprule
Compared model & Difference & 95\% interval & Only proposed model succeeded & Only compared model succeeded & Exact McNemar $p$ \\\midrule
Collision-Aware Teacher & -0.025 & [-0.075, +0.020] & 9 & 14 & 0.4049 \\
Caution-Study Control Model & +0.030 & [-0.005, +0.065] & 10 & 4 & 0.1796 \\
Stock Intelligent Driver Model & +0.040 & [+0.000, +0.080] & 12 & 4 & 0.0768 \\
Baseline Driving Model Trained Without Traffic & +0.075 & [+0.035, +0.120] & 17 & 2 & 0.0007 \\
Baseline Driving Model Trained With Traffic & +0.180 & [+0.125, +0.240] & 38 & 2 & $<0.0001$ \\
Image-Only Driving Model & +0.795 & [+0.740, +0.850] & 159 & 0 & $<0.0001$ \\
\bottomrule\end{tabular}
}
\end{table}

\explainedfigure{fig_per_episode_outcome_strips}{0.97}{Which test episodes each model fails}{Each thin column is one of the 200 test episodes; each row is one model. Green means the model succeeded in that episode, red means it failed. The episodes are sorted so that those failed by no model are on the left and those failed by many models are on the right.}{Most episodes are solved by every model. The failures are concentrated in a small set of hard scenarios on the right, and the models fail many of the same ones. The two baselines also fail a block of episodes that all the other models solve.}{Because failures cluster in particular scenarios, the remaining errors can be studied scenario by scenario instead of treated as random. This is what the failure analysis does, and it is why the comparison is paired (all models see the same scenarios).}{fig:strips}


\explainedfigure{fig_speed_versus_energy_200_episodes}{0.97}{Mean driving speed and estimated energy use of each model}{Each point is one model, placed by its mean speed (horizontal, metres per second) and its estimated traction energy per kilometre (vertical, watt-hours per kilometre) over the 200 clear-weather test episodes. The energy value is the mechanical work at the wheels of a declared virtual vehicle (no regenerative braking), not a battery measurement.}{The two baselines use less energy (85.9 and 90.7 watt-hours per kilometre) and drive faster (8.09 and 8.07 metres per second) than the proposed model (96.2 watt-hours per kilometre at 7.87 metres per second), because they are less careful. The caution-study control is the cheapest (80.4 watt-hours per kilometre).}{Energy comparisons are only meaningful between models with similar safety; a model that crashes early also uses little energy. Differences are valid as comparisons between models, not as absolute battery figures.}{fig:energy}


\subsection{What this result does \emph{not} show}
\begin{itemize}
\item It does not separate the \emph{causes} of the proposed model's gain over the baselines. Three things changed at once relative to the baseline behaviour cloning: (i) dataset aggregation instead of five demonstrations (the main suspect), (ii) the map overlay in the grid, (iii) the split of the vehicle-state inputs. A student with the unmodified grid and all six vehicle-state inputs was evaluated on 10 and 40 episodes only (0.70/0.875, not significantly different from the proposed model's 0.80/0.925);
  \textbf{it has not been evaluated on the 200-episode block}, so the contribution of the overlay at this sample size is untested (Section~\ref{sec:directions}).
\item It does not show anything about dense traffic: no learned model solves it (Section~\ref{sec:failure}).
\item It is a result of the camera-free surrogate (Section~\ref{sec:surrogate}), not of the live camera pipeline.
\end{itemize}
